% Original PDF: source/普通高考/2009/2009山东理.pdf, question 15.
% Same directed topology as the liberal-arts source figure, independently
% checked against the science-paper raster: T>S?; 否 runs through three update
% boxes and loops to the decision, while 是 outputs T and ends.
% solid segments: start-init-test, test-sum-incn-inct, test-output-finish,
%                 inct-loop-test; dashed segments: none.
% labels: 是/否 are placed beside their decision branches; node text is centered
%         with zero paragraph indentation and kept clear of arrows.
\begin{tikzpicture}[
  >=Stealth,
  line width=0.55pt,
  line cap=round,
  line join=round,
  every node/.style={font=\normalsize,anchor=center,align=center,
    execute at begin node={\everypar{}\leftskip=0pt\rightskip=0pt\parindent=0pt
      \hangindent=0pt\hangafter=1\parshape=0}},
  flowbox/.style={draw,minimum height=.70cm,inner xsep=4pt,inner ysep=2pt,
    align=center},
  terminal/.style={flowbox,rounded corners=.18cm,minimum width=1.55cm,
    minimum height=.74cm},
  io/.style={flowbox,shape=trapezium,trapezium left angle=72,
    trapezium right angle=108,minimum height=.80cm},
  decision/.style={flowbox,shape=diamond,aspect=2.65,minimum width=3.45cm,
    minimum height=.95cm,inner sep=1pt},
  flowarrow/.style={->,line width=0.55pt},
]
  \node[terminal] (start) at (0,10.35) {开始};
  \node[flowbox,minimum width=4.25cm] (init) at (0,8.95)
    {$S=0,\ T=0,\ n=0$};
  \node[decision,minimum width=3.70cm] (test) at (0,7.25) {$T>S$};
  \node[flowbox,minimum width=2.90cm] (sum) at (0,5.45)
    {$S=S+5$};
  \node[flowbox,minimum width=2.90cm] (incn) at (0,3.85)
    {$n=n+2$};
  \node[flowbox,minimum width=2.90cm] (inct) at (0,2.25)
    {$T=T+n$};
  \node[io,minimum width=2.45cm] (output) at (4.15,6.05) {输出 $T$};
  \node[terminal] (finish) at (4.15,4.45) {结束};
  \draw[flowarrow] (start.south) -- (init.north);
  \draw[flowarrow] (init.south) -- (test.north);
  \draw[flowarrow] (test.south) -- (sum.north);
  \node[anchor=east] at (-0.18,6.40) {否};
  \draw[flowarrow] (sum.south) -- (incn.north);
  \draw[flowarrow] (incn.south) -- (inct.north);
  \draw[flowarrow] (test.east) -- (4.15,7.25) -- (output.north);
  \node[anchor=west] at (2.10,7.43) {是};
  \draw[flowarrow] (output.south) -- (finish.north);
  \draw[flowarrow] (inct.west) -- (-2.90,2.25) -- (-2.90,7.25) -- (test.west);
  \path[use as bounding box] (-3.35,0.55) rectangle (5.35,10.70);
\end{tikzpicture}
